\documentclass[article]{jss}

\usepackage{orcidlink,thumbpdf,lmodern}
\usepackage{booktabs,tabularx,array}
\usepackage{microtype}
\usepackage[section]{placeins}

\newcolumntype{Y}{>{\raggedright\arraybackslash}X}

\author{Stefanos Balaskas~\orcidlink{0000-0003-2444-9796}\\University of Patras}
\title{\pkg{gp3ml}: Governance-First Predictive Modelling and Validation for Gazepoint Research in \proglang{R}}

\Plainauthor{Stefanos Balaskas}
\Plaintitle{gp3ml: Governance-First Predictive Modelling and Validation for Gazepoint Research in R}
\Shorttitle{\pkg{gp3ml}: Governance-First Predictive Modelling}

\Abstract{
Predictive modelling with eye-tracking data is complicated by repeated participant--stimulus observations, heterogeneous recording quality, flexible preprocessing, and information leakage across resampling partitions. General machine-learning frameworks provide broad modelling functionality, but the validity of a behavioural prediction workflow also depends on the declared outcome, grouping structure, generalization target, tuning boundaries, decision policy, and computational provenance. We introduce \pkg{gp3ml}, an \proglang{R} package for governance-first predictive modelling and validation with Gazepoint-derived and similarly structured behavioural data. The package binds explicit task and feature-role declarations to participant-, stimulus-, trial-, and crossed-group resampling; fold-local preprocessing; nested tuning; retained assessment predictions; calibration and target-aware conformal summaries; threshold and abstention rules; shift and robustness audits; locked analysis plans; portable model artefacts; environment capture; and explicit API contracts. A deterministic synthetic case study uses an observed recording-quality review status to compare row-level, participant-grouped, and participant--stimulus validation designs. The contribution is not a new learner catalogue or automatic model-selection system. It is an auditable governance layer that connects statistical operations to the units and claims a study is intended to support.
}

\Keywords{machine learning, eye tracking, grouped resampling, data leakage, reproducible research, model governance, \proglang{R}}
\Plainkeywords{machine learning, eye tracking, grouped resampling, data leakage, reproducible research, model governance, R}

\Address{
Stefanos Balaskas\\
eGovernment \& eCommerce Lab (Innovation \& Entrepreneurship)\\
Department of Business Administration\\
University of Patras\\
26504 Patras, Greece\\
E-mail: \email{s.balaskas@ac.upatras.gr}\\
ORCID: \url{https://orcid.org/0000-0003-2444-9796}
}

% Draft checkpoint date. Replace with the actual submission date immediately
% before upload to JSS.
\Submitdate{2026-08-08}

\begin{document}

\section{Introduction}
\label{sec:introduction}

Eye-tracking experiments commonly produce multiple observations from each participant,
repeated presentations of stimuli, and measurements nested within trials or recordings.
Such rows are not exchangeable in the ordinary independent-observation sense. A
classifier can exploit stable participant characteristics, recurring stimulus properties,
or duplicated processing decisions when the same clusters occur in both training and
assessment data. Crossed participant--item structures are common in behavioural
research, while scientific claims are usually intended to extend beyond the particular
people, stimuli, or recordings observed in one sample
\citep{Baayen2008,Barr2013,Yarkoni2022}.

A random row split can therefore answer a narrower question than its label suggests. It
may estimate performance for additional rows from already represented participants and
stimuli even when the scientific claim concerns new participants, new stimuli, or both.
The estimate is not intrinsically invalid; it becomes misaligned when interpreted as
evidence for a different target. This distinction motivates the central design rule in
\pkg{gp3ml}: the intended generalization unit is declared before folds are created, and
fold construction preserves that declaration.

Leakage is a second boundary problem. Information can cross the analysis--assessment
boundary through outcome-derived features, post-outcome measurements, identifiers,
globally estimated preprocessing, feature selection, tuning, or threshold optimization
\citep{Kaufman2012,Kapoor2023}. Model selection creates a related boundary: using the
same resampling results to select a model and report its expected performance can bias
the estimate \citep{Varma2006,Cawley2010}. Nested resampling addresses selection bias,
but nesting alone does not determine which participants or stimuli may cross a
partition, how preprocessing is estimated, what uncertainty is reported, or how a
probability becomes a decision.

These concerns motivate a governance view of predictive validity. A defensible workflow
connects the outcome definition, feature provenance, resampling target, tuning rule,
calibration method, threshold policy, robustness evidence, and computational
environment. Performance metrics are necessary, but they do not establish by
themselves that an analysis answered the intended scientific question.

\pkg{gp3ml} \citep{gp3ml2026} implements this connection for Gazepoint-derived and
similarly structured tabular behavioural data. It provides modelling and evaluation
utilities, but its distinctive purpose is to make design and reporting decisions
explicit and auditable. The package does not autonomously identify a ``best'' model,
relax prohibited-use boundaries, or treat predictive associations as evidence of
latent psychological, health, biometric, or other sensitive states.

\section{Predictive validation as a governance problem}
\label{sec:governance-problem}

\subsection{Generalization targets in clustered behavioural data}

The organizing quantity in \pkg{gp3ml} is a machine-readable generalization target.
Holding out participants estimates a different prediction problem from holding out
stimuli or crossed participant--stimulus blocks. Row-level assessment can be suitable
for additional observations from represented clusters, but it should not be reported
as new-participant performance without further justification. Materialized partitions
make these semantics inspectable rather than leaving them implicit in a resampling
label.

\subsection{Leakage and preprocessing boundaries}

Feature manifests distinguish outcomes, predictors, identifiers, grouping columns,
availability stages, outcome-derived features, and preprocessing scope. Fold audits
then test whether the encoded partition respects the declared target. Preprocessing
operations such as imputation, centring, and scaling are fitted on analysis partitions
and applied to assessment rows. Thus the package treats leakage prevention as a
combination of feature provenance and partition-local computation rather than as one
post hoc diagnostic.

\subsection{Selection, uncertainty, and decision boundaries}

Nested resampling separates inner model or hyperparameter selection from outer
performance estimation. Calibration, conformal prediction, threshold selection, and
abstention are treated as distinct stages rather than as interchangeable notions of
``confidence.'' Conformal summaries carry a declared calibration unit and
generalization target, while decision rules record threshold origin, metric direction,
training partition, costs, scientific justification, and optional abstention
boundaries. These objects do not replace substantive judgement; they make the judgement
reviewable.

\section{Related software}
\label{sec:related}

The contribution of \pkg{gp3ml} is intentionally narrower than that of general
machine-learning frameworks. \pkg{caret} unified model training and resampling behind
a common interface \citep{Kuhn2008}. \pkg{mlr3} provides extensible tasks, learners,
resampling, tuning, and benchmarking \citep{Lang2019}; its current task abstraction
also supports a group column role whose observations are kept together during
resampling. The \pkg{tidymodels} ecosystem integrates modelling components with tidy
data principles \citep{KuhnSilge2022}, while \pkg{rsample} provides
\code{group\_vfold\_cv()} for grouped repeated-measures settings and \code{nested\_cv()}
for nested resampling \citep{rsample2026}.

The same statistical mechanisms are not unique to \proglang{R}. In
\proglang{Python}, \pkg{scikit-learn} \citep{Pedregosa2011,scikitLearnDocs2026}
provides \code{GroupKFold} and \code{StratifiedGroupKFold} for non-overlapping groups,
\code{Pipeline} for fitting preprocessing and prediction steps together, and standard
compositions for nested cross-validation. These capabilities are mature and should not
be represented as innovations of \pkg{gp3ml}. Rather, \pkg{gp3ml} adds a
research-governance contract around comparable operations: target labels are retained
with folds and evaluations; feature availability and outcome derivation are auditable;
selection rationales and decision provenance are explicit; and reproducibility evidence
can be validated as package objects.

Other tools address complementary parts of the workflow. \pkg{DALEX} provides
model-agnostic explanation \citep{Biecek2018}; \pkg{targets} manages reproducible
computational pipelines \citep{Landau2021}; \pkg{eyetrackingR} supports preparation and
inferential analysis of eye-tracking measures \citep{DinkFerguson2015}; and RO-Crate
tooling packages research artefacts and metadata \citep{SoilandReyes2022}. Conformal
prediction, probability calibration, and dataset shift also have established
methodological literatures \citep{Vovk2005,Shafer2008,Angelopoulos2023,Guo2017,
QuioneroCandela2009,MorenoTorres2012}. \pkg{gp3ml} integrates conservative forms of
these ideas into one governed predictive workflow without claiming that software alone
resolves their assumptions.

\begin{table}[t]
\centering
\caption{\label{tab:software}Scope-qualified comparison of selected modelling
frameworks. ``External composition'' means that the capability can be assembled by the
user but is not represented as the primary governance contract described here.}
\small
\begin{tabularx}{\textwidth}{@{}lYYY@{}}
\toprule
Software & Grouped validation & Leakage-resistant preprocessing & Research-governance layer\\
\midrule
\pkg{gp3ml} & Target-labelled participant, stimulus, trial, and crossed partitions &
Feature manifests plus fold-local preprocessing and audits &
Purpose, target, decision, robustness, artefact, and API provenance\\
\pkg{mlr3} & Native group role and general resampling infrastructure &
Composable through learner/pipeline infrastructure &
External composition\\
\pkg{rsample}/\pkg{tidymodels} & Native grouped and nested resampling &
Composable recipes/workflows &
External composition\\
\pkg{caret} & User-supplied indices and resampling controls &
Composable through the modelling workflow &
External composition\\
\pkg{scikit-learn} & \code{GroupKFold}, \code{StratifiedGroupKFold}, related splitters &
\code{Pipeline} fits preprocessing and estimators together &
External composition\\
\bottomrule
\end{tabularx}
\end{table}

\subsection{What gp3ml adds}

The package contribution is therefore not grouped cross-validation, pipelines, or
nested tuning in isolation. It is the binding of these operations to explicit research
semantics. In particular, \pkg{gp3ml} records (i) the scientific purpose and observed
outcome, (ii) predictor and grouping roles, (iii) the intended generalization target,
(iv) fold-local preprocessing and leakage evidence, (v) tuning and selection
boundaries, (vi) calibration and uncertainty semantics, (vii) threshold and abstention
provenance, and (viii) release, environment, and artefact evidence.

This added structure has costs. It introduces more ceremony than a minimal modelling
script, its learner catalogue is deliberately smaller than those of general-purpose
frameworks, several function names remain Gazepoint-oriented, and strongly grouped or
crossed designs can reduce the effective assessment sample. These disadvantages are
intentional trade-offs for analyses in which the scientific claim, not algorithmic
breadth, is the central object under review.

\section{Software design and architecture}
\label{sec:architecture}

\pkg{gp3ml} is layered rather than tied to one learner. Task declarations and feature
manifests precede resampling. Preprocessing and model fitting occur inside materialized
partitions. Uncertainty, shift, and decision rules consume retained assessment
predictions. Artefacts, environment captures, and API contracts preserve resulting
evidence. Figure~\ref{fig:architecture} summarizes this architecture.

\begin{figure}[t]
\centering
\includegraphics[width=.90\textwidth]{figures/architecture-figure-1.pdf}
\caption{\label{fig:architecture}Conceptual architecture of \pkg{gp3ml}.
Governance begins before model fitting and continues through release evidence.}
\end{figure}

\begin{table}[t]
\centering
\caption{\label{tab:principles}Core design principles and representative functions.}
\small
\begin{tabularx}{\textwidth}{@{}p{.26\textwidth}Y@{}}
\toprule
Principle & Representative \pkg{gp3ml} mechanisms\\
\midrule
Declared purpose and roles &
\code{declare\_gazepoint\_task()}, \code{declare\_gazepoint\_analysis\_plan()},
\code{validate\_gazepoint\_ml\_roles()}\\
Feature provenance and leakage resistance &
\code{create\_gazepoint\_feature\_manifest()}, \code{audit\_gazepoint\_ml\_leakage()}\\
Target-aligned grouped validation &
\code{create\_gazepoint\_group\_folds()}, \code{create\_gazepoint\_nested\_folds()}\\
Fold-local preprocessing and governed fitting &
\code{fit\_gazepoint\_preprocessor()}, \code{evaluate\_gazepoint\_group\_folds()}\\
Selection and uncertainty &
\code{create\_gazepoint\_tuning\_grid()}, \code{select\_gazepoint\_model()},
\code{fit\_gazepoint\_conformal()}\\
Decision governance &
\code{evaluate\_gazepoint\_thresholds()}, \code{create\_gazepoint\_decision\_rule()},
\code{audit\_gazepoint\_abstention()}\\
Robustness and provenance &
\code{audit\_gazepoint\_dataset\_shift()}, \code{audit\_gazepoint\_model\_robustness()},
\code{capture\_gazepoint\_environment()}\\
Portable evidence and API contracts &
\code{create\_gazepoint\_model\_artifact()}, \code{gp3ml\_api\_contracts()}\\
\bottomrule
\end{tabularx}
\end{table}

At the frozen 0.3.0 release used for this paper, the source contains 127 namespace
exports: 71 stable exports and 56 experimental exports, together with 38 stable public
classes. The release workspace reports 20 source vignettes, 154 Rd topics, and an API
audit with no recorded differences. These counts are release evidence, not claims of
scientific validity.

\section{Using gp3ml}
\label{sec:using}

A minimal governed workflow starts by declaring the outcome and generalization target,
constructing a feature manifest, materializing target-aligned folds, auditing those
folds, and evaluating a learner inside them. The example below uses the same
deterministic recording-quality task as the case study.

\begin{CodeInput}
library("gp3ml")

data <- simulate_gazepoint_governed_data(
  n_participants = 30L, n_stimuli = 8L,
  trials_per_cell = 2L, seed = 20260807L
)
predictors <- c("tracking_ratio", "blink_rate", "gaze_dispersion")

task <- create_gazepoint_synthetic_task(
  data, workflow = "recording_quality",
  generalization_target = "new_participants"
)
manifest <- create_gazepoint_synthetic_manifest(
  task$outcome, predictors
)

folds <- create_gazepoint_group_folds(
  data = data, outcome = task$outcome,
  predictors = predictors, feature_manifest = manifest,
  generalization_target = task$generalization_target,
  participant_id = "participant_id",
  trial_id = "trial_id", stimulus_id = "stimulus_id",
  v = 5L, repeats = 2L, seed = 20260807L
)
audit <- audit_gazepoint_group_folds(folds)
\end{CodeInput}

The audit does not infer whether the research question itself is appropriate. It checks
the encoded contract. For the participant-grouped example, the ten materialized folds
pass their separation checks.

\begin{CodeOutput}
<gazepoint_group_folds_audit>
Overall status: PASS
Audited folds: 10
Non-passing checks: 0
\end{CodeOutput}

Fitting uses the same target and estimates preprocessing within each analysis
partition.

\begin{CodeInput}
evaluation <- evaluate_gazepoint_group_folds(
  folds = folds, task = task, predictors = predictors,
  engine = "glm",
  preprocessor_args = list(
    numeric_imputation = "median",
    center = TRUE, scale = TRUE
  ),
  threshold = 0.5, seed = 20260807L,
  assess_calibration = TRUE, keep_models = TRUE
)

summary <- summarize_gazepoint_resample_performance(evaluation)
\end{CodeInput}

The evaluation object retains the generalization target, fold status, assessment-only
predictions, preprocessing metadata, and metrics. Nested tuning is available when
candidate models or hyperparameters are selected, but it is not used in the main
demonstration because a single logistic regression isolates the effect of the
validation target.

\section{Synthetic case study}
\label{sec:case-study}

\subsection{Data-generating design}

The deterministic synthetic dataset contains 480 rows from 30 participants and eight
stimuli, with two trials per participant--stimulus cell. The observed outcome
\code{quality\_status} records a synthetic \code{pass}/\code{review} recording-quality
status. It is not a psychological construct. Predictors are \code{tracking\_ratio},
\code{blink\_rate}, and \code{gaze\_dispersion}. Five percent of
\code{gaze\_dispersion} values are set to missing with a fixed seed to test whether
imputation remains fold-local.

The locked analysis plan fixes the outcome, predictors, new-participant target,
preprocessing, logistic-regression candidate model, balanced accuracy as the primary
metric, secondary ROC AUC/Brier/log-loss metrics, calibration diagnostics, conformal
calibration semantics, threshold policy, seed strategy, and prohibited
interpretations. The plan is an immutable comparison point for later deviations; it is
not represented as equivalent to a registered preregistration
\citep{Nosek2018}.

\subsection{Validation designs}

Three designs use the same outcome and predictors. The row-level comparator deliberately
allows cluster overlap and is manuscript-specific because \pkg{gp3ml} does not expose
an unlabeled unsafe split as a governed target. Participant-grouped folds keep complete
participants on one side of each split. Crossed participant--stimulus folds prevent
assessment-block participant and stimulus overlap and explicitly retain incompatible
cross-block rows as excluded observations. Figure~\ref{fig:resampling} shows the
different estimands.

\begin{figure}[t]
\centering
\includegraphics[width=.95\textwidth]{figures/resampling-design-figure-1.pdf}
\caption{\label{fig:resampling}Three validation designs applied to the same synthetic
data. The row-level comparator permits represented clusters on both sides; the grouped
designs target new participants or new participant--stimulus blocks.}
\end{figure}

\begin{table}[t]
\centering
\caption{\label{tab:validation-designs}Validation designs and their interpretation.}
\small
\begin{tabularx}{\textwidth}{@{}p{.25\textwidth}YY@{}}
\toprule
Design & Intended estimand & Cluster recurrence\\
\midrule
Naive row-level & Additional rows from represented clusters &
Participant and stimulus recurrence possible\\
Participant grouped & New participants with represented stimulus set &
Participant recurrence prevented\\
Participant--stimulus blocks & New participant--stimulus blocks &
Assessment-block participant and stimulus recurrence prevented\\
\bottomrule
\end{tabularx}
\end{table}

\subsection{Validation-design dependence}

The results are not a benchmark of logistic regression. Their purpose is to show that
the phrase ``model performance'' is conditional on which units may recur across
analysis and assessment data. In this synthetic realization, mean balanced accuracy is
0.500 under all three designs. ROC AUC is 0.603 for the row-level comparator, 0.593 for
participant-grouped validation, and 0.618 for participant--stimulus blocks. Brier scores
are 0.126, 0.128, and 0.123, respectively, while log loss is 0.418, 0.423, and 0.415.

\begin{table}[t]
\centering
\caption{\label{tab:performance}Mean fold-distribution metrics by validation design.
Balanced accuracy is identical in this realization; the values are not intended as a
model benchmark.}
\small
\begin{tabular}{@{}lrrrr@{}}
\toprule
Design & Bal. acc. & ROC AUC & Brier & Log loss\\
\midrule
Naive row-level & 0.500 & 0.603 & 0.126 & 0.418\\
Participant grouped & 0.500 & 0.593 & 0.128 & 0.423\\
Participant--stimulus blocks & 0.500 & 0.618 & 0.123 & 0.415\\
\bottomrule
\end{tabular}
\end{table}

\begin{figure}[t]
\centering
\includegraphics[width=.78\textwidth]{figures/performance-comparison-figure-1.pdf}
\caption{\label{fig:performance}Balanced accuracy by validation design. Error bars are
empirical 95\% fold-distribution quantiles, not independent-observation confidence
intervals.}
\end{figure}

A higher row-level value, had it occurred, would indicate that cluster recurrence made
that prediction problem easier; it would not prove that a grouped estimate was
pessimistic. Similar values instead indicate that these predictors do not benefit
strongly from cluster recurrence in this particular synthetic realization.

\subsection{Calibration and conformal prediction}

Calibration is assessed from retained participant-grouped assessment predictions from
the first repeated-cross-validation pass. The resulting calibration intercept is
-0.755, the slope is 0.547, Brier score is 0.128, log loss is 0.425, and expected
calibration error is 0.021. Figure~\ref{fig:calibration} is descriptive because repeated
trials from the same participant preclude interpreting row-level uncertainty as
participant-level uncertainty.

\begin{figure}[t]
\centering
\includegraphics[width=.72\textwidth]{figures/calibration-figure-1.pdf}
\caption{\label{fig:calibration}Reliability diagram for participant-grouped assessment
predictions. Point size reflects observations per probability bin.}
\end{figure}

For conformal prediction, retained predictions are divided into participant-disjoint
policy-development and policy-audit partitions. Conformity scores are estimated from 15
development participants, aggregated to the participant calibration unit, and evaluated
on 15 held-out participants. At nominal coverage 0.90, held-out row coverage and the
proportion of participants with all rows covered are both 1.000; mean prediction-set
size is 1.988 classes. These results are an executable illustration, not a claim of
distribution-free coverage under arbitrary dependence
\citep{Vovk2005,Shafer2008,Angelopoulos2023}.

\begin{table}[t]
\centering
\caption{\label{tab:calibration}Calibration and participant-disjoint conformal
summaries.}
\small
\begin{tabular}{@{}lr@{}}
\toprule
Quantity & Value\\
\midrule
Calibration intercept & -0.755\\
Calibration slope & 0.547\\
Expected calibration error & 0.021\\
Nominal conformal coverage & 0.900\\
Held-out row coverage & 1.000\\
Held-out participant all-rows coverage & 1.000\\
Mean prediction-set size & 1.988\\
\bottomrule
\end{tabular}
\end{table}

\subsection{Decision policy}

A probability estimate is not itself a decision rule. Candidate thresholds are
evaluated on one participant subset and the resulting rule is audited on disjoint
participants. The selected threshold is 0.20 and is recorded as training-derived.
An illustrative abstention interval of 0.15--0.25 is enabled as a reporting device
rather than an operational recommendation. On the 240 held-out rows in this
realization, coverage is 0.621, the abstention rate is 0.379, and error among
covered cases is 0.221.

\begin{figure}[t]
\centering
\includegraphics[width=.90\textwidth]{figures/threshold-abstention-figure-1.pdf}
\caption{\label{fig:decision}Threshold development on one participant subset and
abstention auditing on disjoint held-out participants.}
\end{figure}

\subsection{Dataset shift and robustness}

A second synthetic dataset is perturbed transparently: \code{tracking\_ratio} is shifted
downward, \code{gaze\_dispersion} is scaled, and its missingness increases. The shift
audit reports standardized differences of -1.974 for \code{tracking\_ratio}, 0.221 for
\code{blink\_rate}, and 1.503 for \code{gaze\_dispersion}. The first and third findings
are flagged \code{fail}; \code{blink\_rate} is flagged \code{review}. Missingness is
audited separately so that one scalar does not conflate distributional and missing-data
changes.

\begin{figure}[t]
\centering
\includegraphics[width=.78\textwidth]{figures/shift-figure-1.pdf}
\caption{\label{fig:shift}Predictor-distribution shift between development and
perturbed synthetic external data.}
\end{figure}

The aggregate robustness audit is \code{fail}: seed sensitivity passes, whereas the
threshold-stability diagnostic fails under the encoded acceptance rule. A failure here
is intentionally retained rather than normalized away. These diagnostics do not
replace external validation and do not identify a causal source of drift.

\section{Discussion}
\label{sec:discussion}

\pkg{gp3ml} treats governance as part of predictive validity. This does not imply that
every analysis needs the same split, learner, metric, or decision rule. It means that
these choices define the claim being made and should therefore be explicit,
machine-readable, and reviewable.

The package complements rather than replaces broader modelling ecosystems. Researchers
requiring extensive learner catalogues, sophisticated search spaces, distributed
benchmarking, or mature cross-language tooling may prefer \pkg{mlr3},
\pkg{tidymodels}, \pkg{scikit-learn}, or related infrastructure and use
\pkg{gp3ml}-style governance as an additional layer. Conversely, a transparent
scientific analysis may require nothing more complex than the generalized linear model
used here. The contribution is workflow integrity rather than algorithmic complexity.

The device-specific scope is similarly bounded. Although the package was developed
around Gazepoint-derived research data, its grouping and governance abstractions are
not tied to a file format. Tabular behavioural data with declared outcomes, predictors,
participants, stimuli, trials, or other clusters can use the same logic. Device-specific
signal processing should occur upstream and its transformations should be represented
in provenance or handoff metadata.

\subsection{Limitations and disadvantages}

Software governance cannot repair a weak research design, invalid outcome, inadequate
measurement, or unrepresentative sample. A passing leakage audit means that encoded
checks passed; it does not prove that every temporal, causal, or procedural dependency
has been identified.

Grouped validation can expose an effective sample size far smaller than the row count.
With few participants or stimuli, folds may be unstable or lack an outcome class.
Crossed designs may explicitly exclude many incompatible rows, trading nominal sample
size for a more stringent generalization target. \pkg{gp3ml} retains warnings and
failed folds, but the researcher must respond through redesign, simpler claims, or
additional data.

The workflow also adds object and reporting overhead relative to a concise
general-purpose machine-learning script. Its governed engine catalogue is intentionally
limited, and optional engines have heterogeneous installation requirements. The package
does not provide unrestricted AutoML, distributed benchmarking, or an exhaustive model
search interface.

The case study is synthetic. It demonstrates execution and interpretation but does not
establish performance across real eye-tracking devices, tasks, populations, or recording
conditions. Broader independent benchmarking is future work. Conformal summaries depend
on exchangeability assumptions and calibration-unit definitions; grouped aggregation is
conservative but does not provide distribution-free guarantees under arbitrary
dependence. Dataset-shift audits describe measured changes but do not guarantee
performance degradation.

\section{Reproducibility and software availability}
\label{sec:reproducibility}

The manuscript is paired with a single commented replication script,
\code{code.R}, and its spun output, \code{code.html}. The script recreates all six
figures used in the article, exports the principal result tables, and finishes with
\code{sessionInfo()}. The synthetic generator uses a fixed seed. This follows the JSS
requirement that all displayed figures, tables, and outputs be reproducible and the
journal's recommendation that an \proglang{R} submission provide
\code{code.html} generated from \code{code.R}.

\begin{sloppypar}
Version 0.3.0 is archived as GitHub tag \code{v0.3.0}. The repository is
\url{https://github.com/stefanosbalaskas/gp3ml} and the documentation site is
\url{https://stefanosbalaskas.github.io/gp3ml/}. The Zenodo concept DOI is
\url{https://doi.org/10.5281/zenodo.21487272}. At this draft checkpoint, 0.3.0 is a
GitHub source release; it is not represented here as a CRAN 0.3.0 release. The package
uses an MIT licence, which is GPL-compatible for the JSS software-licensing
requirement.
\end{sloppypar}

The frozen 0.3.0 source archive used for release validation has SHA-256
\texttt{04c5fdb0\allowbreak{}17c63ecb\allowbreak{}f6f5e85c\allowbreak{}29fff649\allowbreak{}2880201a\allowbreak{}2f2aa5c3\allowbreak{}c1d3fe99\allowbreak{}8e60bae6}.
The release checks completed with zero errors, zero warnings, and zero notes. These are
software-release facts, not evidence that every scientific application of the package
is valid. Reproducibility and reporting practices follow the broader emphasis on
executable artefacts and transparent machine-learning workflows
\citep{Pineau2021}.

\section{Conclusions}

\pkg{gp3ml} provides a governance-first layer for predictive modelling with clustered
behavioural data in \proglang{R}. Its contribution is not a new cross-validation
algorithm or learner catalogue. Instead, it binds familiar statistical operations to
declared scientific purposes, feature provenance, generalization targets, selection
boundaries, uncertainty semantics, decision provenance, robustness evidence, and
computational artefacts. The case study shows how the same data and model can be
evaluated under different target semantics while keeping those distinctions explicit
in the resulting objects. This makes predictive claims easier to inspect, reproduce,
and review without replacing the scientific judgement on which those claims depend.

\bibliography{refs}

\end{document}
